\section{Reproducibility and mathematical scope}
\label{sec:verification-release}
The numerical claims in this paper arise from three
different forms of calculation, and their logical roles
should be kept separate. The fifth-iterate calculation
for the genuine theta kernel uses outward rational
integer intervals, a Machin-series enclosure for
$\pi$, and explicit control of the omitted theta
summands. Its certificate is reproduced by
\path{certificates/xi_fifth_iterate_exact.py}.
The Gaussian-polynomial identities, the
large-frequency truncation coefficients and
the rational bounds on the exponential theta
tail can be checked directly from the
recurrences given in the proofs.

The radial monotone-curvature property of the classical
theta potential was previously proved by Newman in 1991
using a GHS inequality \cite{NewmanGHS}. The certificate
provided here is therefore an alternative quantitative
verification, not the discovery of that property.
The radial-curvature certificate is a
computer-assisted analytic result.
Its proof separates a compact interval,
certified by 500 outward-rounded rational
cells, from a distant region controlled
by an explicit infinite-tail majorant.
The implementation and the relevant
polynomial identities appear in
\path{certificates/xi_global_curvature_exact.py}.
Although the arithmetic checks are exact,
an independent mathematical review of
the complete tail comparison and the
placement of this theorem within earlier
theta-kernel literature remains appropriate.

The four finite zero-count theorems certify
distinct disks. The first three use a two-mode
Fourier approximation; the fourth requires
three modes and tighter local derivative bounds.
Their proofs reduce to enclosures of
three Taylor coefficients and one
positive remainder integral. Every
coefficient is computed at two
different composite-Simpson resolutions;
the stated fourth-derivative bounds
and the omitted $r>2$ integral tails
are handled analytically.
The integration nodes use
\texttt{mpmath.iv} version 1.3.0 for
interval enclosures of $\pi$,
exponentials and trigonometric values,
which are subsequently converted to
exact rational endpoints.
Thus the computational trust boundary
includes that library's transcendental
interval implementation as well as
the elementary analytic estimates.
The scripts, available in the
\texttt{certificates} directory,
may be run independently of the
LaTeX build. The corresponding tests
and their results are recorded
in the public repository's
\texttt{RESULTS.md} file.

The Lean 4.24.0/Mathlib development
certifies exact algebraic identities,
moment-positivity statements, rational
comparison gates, and the functional
symmetry of an entire form of
the classical completed zeta function.
For these small algebraic theorems,
the recorded axiom inspection reports
the ordinary Mathlib axioms
\texttt{propext}, \texttt{Classical.choice},
and \texttt{Quot.sound}.
The large rational-jet test instead
uses \texttt{native\_decide}, which
also relies on compiler-trust axioms;
it is not a small-kernel-only check.
Neither the interval-Simpson
calculations nor the full
Jacobi transformation, Fourier
integrals and Rouch\'e transfers
have been formalized end to end
in Lean.

The existence and simplicity of
the four zeros certified here
are classical results. Indeed,
Platt and Trudgian verified the
critical-line location and
simplicity of all nontrivial
zeta zeros up to height
$3\cdot10^{12}$ \cite{PlattTrudgian}.
The present certificates therefore
make no claim of a new zeta-zero
location or an improvement in
verified height. Their value is
to test the particular theta-Fourier
approximation and its explicit
Rouch\'e error bounds.

The open difficulty remains global.
The explicit theta truncation estimate
can be made arbitrarily small on the
complex strip, with an adaptive
number of modes growing as
$O(\sqrt{\log R})$ for prescribed
polynomial accuracy at height $R$.
What is unavailable is a
height-uniform strict lower bound
on the approximating transform
over contours that would exclude
off-critical zeros. Nor do the
finite Laguerre tests supply that
missing lower bound, as the
counterexamples in this paper show.
Consequently neither the four
bounded zero counts nor the
formal algebraic results prove
the Riemann hypothesis.

The complete source, reproducibility
instructions and verification record
are maintained at
\url{https://github.com/tonynewton79-web/RiemannHypothesis}.
